% Working paper, refreshed 8 October 2026 from writing commit 5fe730d. Main source for subsequent editing.
% Evidence and revision history: README.md and the preserved local revisions.
% Based on the supplied AAMAS 2027 template; class and layout unmodified.
\documentclass[sigconf,anonymous]{aamas}
\usepackage{balance}
\usepackage{tikz}
\usetikzlibrary{arrows.meta,positioning,calc,fit,backgrounds}
\usepackage{placeins}
\usepackage{float}
% Prevent flush-bottom stretching from opening intermittent section gaps.
\raggedbottom
% Allow paragraph line breaks to absorb long compounds without changing margins.
\emergencystretch=2em
\makeatletter
\gdef\@copyrightpermission{
  \begin{minipage}{0.2\columnwidth}
   \href{https://creativecommons.org/licenses/by/4.0/}{\includegraphics[width=0.90\textwidth]{by}}
  \end{minipage}\hfill
  \begin{minipage}{0.8\columnwidth}
   \href{https://creativecommons.org/licenses/by/4.0/}{This work is licensed under a Creative Commons Attribution International 4.0 License.}
  \end{minipage}
  \vspace{5pt}
}
\makeatother
\setcopyright{ifaamas}
\acmConference[AAMAS 2027]{Proc.\@ of the 26th International Conference
on Autonomous Agents and Multiagent Systems (AAMAS 2027)}{3--7 May 2027}{Hanoi, Vietnam}{M.~Baldoni, F.~Fang, W.~Yeoh, N.~Yorke-Smith (eds.)}
\copyrightyear{2027}
\acmYear{2027}
\acmDOI{}
\acmISBN{}
\acmSubmissionID{2613}
\submissionType{Research Paper Track}
\title{When to Intervene: Comparing Constant and Physiology-Triggered Robotic Assistance}
% Anonymous review copy; the class displays submission ID2613.
\author{Anonymous Author(s)}
\keywords{Human--robot interaction, adaptive assistance, intervention timing, physiological sensing, Wizard-of-Oz, tangram}
\citestyle{acmnumeric}

\begin{abstract}
In human--robot interaction (HRI), deciding when to provide assistance requires understanding how intervention strategies affect both task performance and the user experience. Previous work compared physiology-triggered robotic assistance with periodic support in surgical training. We extend this work to physical spatial problem solving through a within-participant experiment comparing no assistance, constant assistance every 30 seconds, and adaptive assistance triggered by heart-rate signals. Twenty-four participants solve nine physical tangram puzzles across three condition blocks, with balanced condition orders and randomly assigned puzzles. Tangrams combine spatial reasoning with physical manipulation, providing a simple task for studying informational and physical assistance. Participants receive on-screen hints and robotic piece retrieval through a Wizard-of-Oz setup. We evaluate correct-piece progress, completion, round duration, adapted seven-point NASA-TLX workload, intervention experience, frustration, and final trust. Completion rates are 15.3\% without assistance, 59.7\% with constant assistance, and 50.0\% with adaptive assistance. Both assisted conditions improve progress, completion, and workload relative to control (Holm-adjusted $p<.001$). Constant assistance has the highest observed completion and shortest duration; adaptive assistance has descriptively higher helpfulness and timing ratings and lower frustration, but no assisted-policy contrast survives multiplicity correction. These findings support assistance in physical spatial problem solving while leaving the choice between periodic and physiology-triggered assistance unresolved.
\end{abstract}

\begin{document}
\maketitle


\section{Introduction}

A pause during a physical puzzle can reflect productive mental rotation or uncertainty about the next move. A helpful intervention in the latter case may interrupt reasoning in the former. Assistance timing therefore matters alongside assistance content.

Recent systems learn assistance type, timing, and confidence \cite{andriella2025}, estimate action completion for proactive coordination \cite{delazzari2025}, or adapt task-level timing and role allocation \cite{karbouj2026}. Physiological information offers another timing cue: Yang et al. compared workload-adaptive robotic suction with periodic support in surgical training \cite{yang2024}. We extend that comparison to physical spatial problem solving, where text hints guide interpretation and an arm retrieves relevant pieces.

The usefulness of an intervention depends on both its task benefit and its cost to the recipient. Assistance-appropriateness judgments reflect this balance \cite{ramnauth2026}, and a proactive escape-room robot elicited more interaction without significantly improving overall performance \cite{vitry2026}. Evaluating a timing policy therefore requires task outcomes alongside the experience of receiving help, rather than treating more interaction as an improvement in itself.

Physiological triggering also needs task-specific evaluation. Reviews report heterogeneous workload measures and mixed cardiac findings in human--robot collaboration \cite{pereira2025}, while cardiac measurement guidelines distinguish heart rate from heart-rate variability and emphasize signal quality and physical activity \cite{quigley2024}. We use a short-term heart-rate rise as an intervention cue, not as a direct measure of frustration, and evaluate its consequences through task performance and participant ratings.

Tangrams combine spatial reasoning and manipulation using common materials, observable outcomes, and repeated targets; related HRI and assembly research uses the same task family \cite{tabatabaei2025,melo2026}. Each participant attempts three puzzles in each of three conditions. We ask \textbf{RQ1:} How do the conditions differ in correct-piece progress, completion, and round duration? \textbf{RQ2:} How do they differ in reported workload and frustration, and how do the assisted policies compare in intervention experience?

Our contributions are:
\begin{enumerate}
\item A controlled extension of periodic versus physiology-triggered robotic assistance to physical spatial problem solving, with an unassisted comparator and common hints and piece retrieval.
\item Participant-level evidence that both assistance policies support correct-piece progress, completion, and reported workload relative to no assistance, without establishing superiority of either assisted policy.
\item An assessment of intervention experience and final-survey feedback on hint frequency and robotic piece retrieval.
\end{enumerate}

\section{Related Work}

\subsection{Assistance timing and proactive coordination}
Beyond task-state adaptation, gaze can indicate readiness for joint action: Lavit Nicora et al. analyze assembly gaze and pilot gaze-based initiation \cite{lavitnicora2024}. Tanneberg et al.'s LLM-based Attentive Support uses scene and dialogue information to decide whether to assist or remain silent, evaluated through constructed scenarios and a robot demonstration rather than a controlled participant comparison \cite{tanneberg2024}. A mentalising/Q-learning memory-game study reports better performance and acceptance, but explanation detail differs between conditions and timing remains fixed \cite{andriella2025mentalising}. These designs do not isolate physiological triggering against periodic assistance.

Timing also affects feedback and interruption handling. Reminders before robot behavior changes elicit faster, more immediate feedback in Space Invaders, but not a significant whole-game timing effect \cite{candon2023}. Exploratory analyses of user-initiated interruptions associate unsuccessful handling with lower perceived inclusion and discussion satisfaction \cite{cao2025}. Initiating assistance, soliciting feedback, and responding to interruptions are distinct decisions; our comparison concerns assistance initiation.

\subsection{Physiological information in assistance}

Yang et al.'s surgical system uses EEG and eye tracking to inform adaptive suction, with a periodic comparator selected to approximate earlier observed assistance frequency \cite{yang2024}. Teo et al. use individualized physiological workload markers to trigger aid during robot supervision, imposing aid later if no trigger occurs \cite{teo2018}. Both inform our timing comparison, although the tasks, sensors, and assistance differ.

Physiological measurement does not necessarily drive online adaptation. Hostettler et al. adapt behavior to user distance while measuring pupil responses \cite{hostettler2025}; Ojsteršek et al. personalize robot parameters using a skills test and analyze ECG afterward \cite{ojstersek2024}. Korivand et al. develop task-load prediction and Q-learning adjustment but report that recorded wristband data could not be integrated directly in real time \cite{korivand2024}. Prajod et al. evaluate an ECG/facial-expression challenge classifier offline; their adaptive assembly condition uses Wizard-of-Oz judgments of progress rather than physiological triggers \cite{prajod2024}.

Beyond robotic execution, GuideAI combines cardiac, gaze, and behavioral information to adapt learning content, pacing, and feedback \cite{shukla2026}. A single-surgeon simulation study identifies subjective workload and mean heart rate among influential task-performance prediction features \cite{wei2025}. Generalizing such signals remains challenging: Bhagat Smith and Adams review portability, model complexity, and adaptability across unknown tasks \cite{bhagatsmith2026}, and Capponi et al. find no clear RMSSD pattern across assembly configurations \cite{capponi2024}. Our study evaluates an HR-rise timing policy rather than a workload-prediction model.

\subsection{Physical tasks and participant experience}

Tangram studies examine gaze around robot failures \cite{tabatabaei2025}, task demand through contour visibility in SensCogAR \cite{melo2026}, and robotic reasoning and manipulation from silhouette targets in MRChaos \cite{zhao2025}. A preliminary three-participant assembly study compares manual, collaborative, and guided conditions, using EEG to assess workload \cite{caiazzo2024}. These settings connect spatial interpretation with physical manipulation.

Cavicchi et al. caution that social cues can distract as well as support performance and recommend task delegation for cognitive offloading \cite{cavicchi2025}. In assembly, human-led collaboration lowers five workload dimensions, while slower robot pacing lowers mental and temporal demand; pacing changes action onset, not movement speed \cite{vandijk2023}. Responses also vary with population: older adults report greater workload with robot than human guidance in a modified Trail Making Test, whereas the younger-adult difference is not significant \cite{varrasi2026}. Workload definitions differ across tasks; Smit et al.'s simulated order-picking model represents workload fairness through lifted mass \cite{smit2024}.

Assistance format matters too. In a medical-training simulation, a Wizard-of-Oz robotic crash cart's verbal object-search guidance with visual reminders yields lower workload and higher usefulness and ease-of-use ratings than reversed cues or no feedback \cite{tanjim2025}. We keep the available hint and retrieval formats common across assisted policies and compare their timing.

\subsection{Wizard-of-Oz evaluation}

Wizard-of-Oz methods enable human-operated delivery of robot interactions. Thunberg et al.'s workshop proposal emphasizes the practical, ethical, and methodological tensions of the wizard's role \cite{thunberg2026}. Interviews with six HRI researchers identify operator response processing, robot delays, unpredictable participants, and control precision as challenges \cite{bejarano2024}. Our apparatus separates the timing trigger from the researcher's selection and delivery of task-relevant assistance.

\section{Formative Assessment of Tangram Solving}

\subsection{Assessment procedure and task experience}
A formative assessment collected 50 substantive responses: 38 undergraduates, nine graduate students, and three faculty members. Ages ranged from 18 to 30 years ($M=21.96$, $SD=3.57$); 25 respondents identified as men, 20 as women, and five preferred not to disclose gender. A Google Form shared through college mailing groups linked to an online square tangram with a timer. Respondents attempted it and reported their experience and expected assistance preferences. Some also participated in the robot experiment; the samples are not treated as independent replications.

Twenty-nine had previously solved tangrams, including four regular solvers; ten had heard of them without solving them and eleven were unfamiliar. Twenty-eight reported complete solutions (56\%), 18 partial solutions (36\%), and four no solution (8\%). Thirty-seven (74\%) felt stuck at least once. Twenty-six selected a solving time below one minute, 16 selected one to five minutes, three selected five to ten minutes, and five selected not attempting or solving the puzzle. Completion and time are self-reported categories, not instrumented performance measures.

\subsection{Assistance expectations and task rationale}
Respondents independently rated hypothetical assistance from 1 (most helpful) to 5 (least helpful), rather than ranking the options. Mean ratings were 2.58 ($SD=1.34$) for a simpler puzzle, 2.76 ($SD=1.22$) for a text hint, 2.80 ($SD=1.14$) for a video tutorial, and 2.94 ($SD=1.27$) for a robot showing piece placement (Figure~\ref{fig:survey}). Responses requested first-piece hints, visual placement guidance, and help at difficult moments. These descriptive expectations motivate the task and assistance modalities; they do not demonstrate robot efficacy.

\begin{figure}[H]
\centering
\includegraphics[width=\columnwidth]{figures/formative-helpfulness.pdf}
\caption{Expected helpfulness in the formative assessment ($N=50$). Dots show means and bars show SD; lower 1--5 ratings indicate greater expected helpfulness.}
\label{fig:survey}
\Description{Four independently rated assistance options have means between 2.58 and 2.94 on a one-to-five scale.}
\end{figure}
Each experimental tangram requires all seven pieces to reproduce a target shape: two large triangles, one medium triangle, two small triangles, a square, and a parallelogram. The seven colored pieces together form a $10\times10$-inch square. We use nine colored-piece targets, with solutions available to the researcher.

Participants interpret the target, rotate and place pieces, and revise arrangements without a prescribed sequence. Hints address orientation or adjacency; retrieval brings a relevant piece to the workspace. Common materials support repeated comparisons across targets.

\section{Study Method: Assistance System and Experimental Protocol}

\subsection{Architecture and apparatus}

Our system coordinates three browser interfaces through a local server (Figure~\ref{fig:architecture}). The researcher dashboard displays the puzzle solution, workspace camera view, timer, physiological feedback, and assistance controls. The participant laptop presents instructions, text hints, the timer, and questionnaires. The robot-operator display identifies the piece-specific program to execute. The server synchronizes the interfaces and records task outcomes, interventions, and questionnaire responses.

\begin{figure*}[t]
\centering
\begingroup
% Keep labels readable after scaling the supplied standalone diagram.
\renewcommand{\footnotesize}{\fontsize{10}{11}\selectfont}
\renewcommand{\scriptsize}{\fontsize{8.5}{9.5}\selectfont}
\renewcommand{\small}{\fontsize{10}{11}\selectfont}
\resizebox{.86\textwidth}{!}{\input{figures/study-system-diagram.tikz.tex}}
\endgroup
\caption{Assistance architecture. Cardiac measurements pass through signal-quality checks and a heart-rate-rise detector. The assigned condition determines assistance timing; task context guides the researcher’s selection of hints and robot actions. The local server synchronizes interfaces and records study measures.}
\label{fig:architecture}
\Description{Four modules show physiological processing, assistance timing, Wizard-of-Oz delivery and participant interaction. Current and reference heart-rate windows feed an arousal flag. Constant and adaptive triggers lead to hints or piece retrieval, coordinated through a local server.}
\end{figure*}

The setup combines an Orangewood OWL63 robotic arm, a Maxim H Band wrist sensor, and a Logitech C270 workspace camera. An opaque partition separates the participant from two researchers, who administer the task and operate the arm (Figure~\ref{fig:setup}). The camera is mounted on top of the partition and covers the participant table. The participant faces a workspace with a puzzle assembly area in front of and to their left, a drop spot to their right, and a laptop for hints and questionnaires. Unused pieces occupy fixed, labeled pickup positions. The arm retrieves a selected piece using one of seven piece-specific programs and delivers it to the drop spot; the participant places it in the target arrangement. Preset text hints describe piece location, orientation, or adjacency, providing consistent assistance content across participants. Hints are accompanied by an audible notification. The researcher selects either a hint alone or a robot action paired with a hint, based on the current arrangement.

\begin{figure*}[t]
\centering
\begingroup
\begin{minipage}[c]{.70\textwidth}
\centering
\resizebox{\linewidth}{!}{\input{figures/physical-setup-topdown.tikz.tex}}
\end{minipage}\hfill
\begin{minipage}[c]{.28\textwidth}
\centering
% Presentation-only crop; original photograph pixels are retained in the source.
\includegraphics[width=\linewidth,trim=180bp 150bp 420bp 0bp,clip]{figures/setup-wide-photo.jpg}\par\smallskip
\includegraphics[width=\linewidth]{figures/setup-topdown-photo.png}\par\smallskip
\includegraphics[width=\linewidth]{figures/operator-dashboard.png}
\end{minipage}
\endgroup
\caption{Wizard-of-Oz setup: schematic (left); robot and participant workspace, overhead view, and operator dashboard (right, top to bottom). The partition separates researchers from the participant, with a camera observing the table. Dashboard values illustrate the interface, not aggregate results. Schematic not to scale.}
\label{fig:setup}
\Description{The schematic places researchers behind a partition and the participant at a table with pickup positions, a drop spot, an assembly area and a laptop. Three images are vertically centered to its right: a rear view of the participant with the robot arm, partition-mounted camera and laptop; an overhead view of tangram manipulation; and the operator dashboard with workspace, solution preview, timer, physiological panel and assistance controls.}
\end{figure*}

% Allow apparatus floats to flow without shortening the following text column.

\subsection{Physiological monitoring}

After receiving the study information, participants wore the wrist sensor during a three-minute pre-task baseline window. The sensor streams cardiac measurements over Bluetooth to a Python collector. To detect a short-term rise in heart rate, the collector compares median heart rate over the latest 15 seconds with median heart rate over the preceding 60 seconds. The arousal flag is raised when the difference is at least 5 beats per minute. A calibrated baseline provides the reference while the rolling reference is being established. Signal-quality checks exclude out-of-range readings and readings with poor sensor contact, and the detector requires sufficient usable samples before raising a flag.

The dashboard displays the flag alongside the workspace view. In the adaptive condition, the flag initiates assistance, while the puzzle arrangement determines the content of the hint or robot action. The flag is a heart-rate-rise cue for assistance timing; participant workload and frustration are assessed through questionnaires.

\subsection{Assistance conditions}

\textbf{No assistance} provides no text hints or robot interventions. \textbf{Constant assistance} offers help every 30 seconds. \textbf{Physiology-informed adaptive assistance} offers help when the heart-rate-rise flag is raised; repeated flags permit further help, with a 15-second wait for rapid or sustained recurrence. In both assisted conditions, the researcher selects a preset hint alone or a robot intervention paired with a hint according to the current arrangement. Thus, the assigned policy determines timing, while task state determines content.

\subsection{Participants and procedure}

Twenty-four participants volunteered through forms shared in college groups and received coupons. Most were undergraduates, with graduate students/researchers and faculty also participating. Ages ranged from 20 to 27 years ($M=21.96$, $SD=2.01$). Fourteen were men and ten were women (58.3\% and 41.7\%, respectively). Each attempted three blocks of three puzzles, with one condition per block. All six condition orders were represented by four participants each (Figure~\ref{fig:order}); puzzles were randomly assigned to blocks. Consent and background information preceded the task. A solution used all seven pieces to resemble the target; unsolved attempts had a nominal five-minute limit. Questionnaires followed each puzzle, with five-minute breaks between blocks and a final survey after the ninth puzzle.

\begin{figure}[!htbp]
\centering
\begingroup
\fontsize{12}{13}\selectfont
\renewcommand{\small}{\fontsize{11}{12}\selectfont}
\resizebox{\columnwidth}{!}{\input{figures/condition-order-tree.tikz.tex}}
\endgroup
\caption{Balanced condition exposure: four participants per order. C denotes control, K constant, and A adaptive; sequences read left to right.}
\label{fig:order}
\Description{Twenty-four participants divide equally among the six permutations of control, constant and adaptive assistance.}
\end{figure}

\subsection{Measures and analysis}

We measure binary completion, attempt duration, and correct-piece counts recorded in researcher notes and photographs of final arrangements. Counts range from 0 to 7; solved rounds receive 7. Where multiple orientations are admissible, we retain the orientation yielding the highest correct-piece count. Attempt duration includes solved and unsolved puzzles.

Round-specific questionnaires assess six adapted NASA-TLX items \cite{hart2006} on seven-point scales (1--7): mental, physical, and temporal demand, performance, effort, and frustration. Performance runs from failure to perfect and is reversed as $8-P$. We average the six aligned burden ratings without weighting and rescale the mean $B$ as $(B-1)100/6$. This is an adapted unweighted 0--100 workload score, not the original weighted NASA-TLX instrument. Assisted rounds also assess helpfulness, timing, seamlessness, and frustration relief. Final helpfulness, efficacy, trust, and reliance ratings summarize the whole study, not individual conditions.

All 24 participants contribute task outcomes, durations, and round-specific ratings, with 72 puzzle attempts per condition. We compare conditions using two-sided paired $t$ tests on each participant's three-round condition averages and report mean differences, 95\% confidence intervals, and Cohen's $d_z$. To account for multiple comparisons, we apply Holm correction within three groups: completion and overall workload (six comparisons), correct-piece counts (three), and secondary measures (25 across duration, workload dimensions, and assistance ratings). Statistical significance is assessed at adjusted $p<.05$.

\section{Findings}

\subsection{Task performance}
Both assisted conditions improved task performance relative to control (Table~\ref{tab:performance}; Figure~\ref{fig:performance}). Across 216 attempts, 90 puzzles were completed: 15.3\% in control, 59.7\% with constant assistance, and 50.0\% with adaptive assistance.

\begin{table}[H]
\centering
\caption{Performance ($n=24$): participant-condition means (SDs) for pieces and duration; completion totals across 72 attempts per condition.}
\label{tab:performance}
\begin{tabular}{lrrr}
\toprule
Measure & Control & Constant & Adaptive \\
\midrule
Pieces / 7 & 3.43 (1.40) & \textbf{5.68} (1.05) & 5.43 (1.13) \\
Completed / 72 & 11 & \textbf{43} & 36 \\
Duration (s) & 292.4 (22.5) & \textbf{248.2} (53.3) & 265.4 (29.1) \\
\bottomrule
\end{tabular}
\end{table}

Relative to control, constant and adaptive assistance increased correct-piece counts by 2.25 and 2.00 pieces and completion by 44.4 and 34.7 percentage points, respectively (all Holm-adjusted $p<.001$). Constant had higher observed progress and completion than adaptive, but these policy contrasts were not significant ($p=.273$ and $p=.296$, respectively).

Mean attempt duration was 44.2 seconds shorter under constant assistance and 27.0 seconds shorter under adaptive assistance than control (both secondary-adjusted $p=.006$). The adaptive--constant contrast was not significant ($p=1.000$). These durations include unsuccessful attempts and occasional overruns beyond the nominal five-minute limit, rather than measuring time to successful completion alone.

\begin{figure*}[tb]
\centering
\includegraphics[width=\textwidth]{figures/performance-tests.pdf}
\caption{Paired contrasts ($n=24$): mean differences and pointwise 95\% CIs, with $t$ statistics and Holm-adjusted $p$ values (three-test piece or six-test core family). K: constant; C: control; A: adaptive. Blue indicates adjusted $p<.05$; gray, $p\geq.05$.}
\label{fig:performance}
\Description{Both assisted conditions outperform control on piece counts and completion and have lower workload. Adaptive-constant confidence intervals cross zero for all three measures.}
\end{figure*}

\subsection{Subjective workload}
Adapted workload averaged 62.46 in control, 49.50 with constant assistance, and 46.06 with adaptive assistance (Table~\ref{tab:workload}; Figure~\ref{fig:performance}). Reductions against control were 12.96 points for constant (95\% CI [8.52, 17.40], $t(23)=-6.04$, $d_z=-1.23$) and 16.40 for adaptive (CI [10.84, 21.95], $t(23)=-6.10$, $d_z=-1.25$), both core-adjusted $p<.001$. The adaptive--constant difference was not significant ($p=.296$).

\begin{table}[H]
\centering
\caption{Workload means (SDs), $n=24$. Items use 1--7: higher performance is better; higher demand, effort, and frustration indicate greater burden. Overall workload uses 0--100.}
\label{tab:workload}
\setlength{\tabcolsep}{3pt}
\begin{tabular}{lrrr}
\toprule
Measure & Control & Constant & Adaptive \\
\midrule
Mental demand & 5.71 (0.94) & 4.82 (0.98) & \textbf{4.63} (1.46) \\
Physical demand & 3.65 (1.68) & 3.28 (1.37) & \textbf{3.06} (1.60) \\
Temporal demand & 3.92 (1.38) & 3.63 (1.25) & \textbf{3.56} (1.45) \\
Performance & 2.93 (1.39) & 4.54 (1.47) & \textbf{4.94} (1.23) \\
Effort & 5.44 (0.69) & \textbf{4.71} (0.95) & 4.78 (1.19) \\
Frustration & 4.69 (1.32) & 3.93 (1.23) & \textbf{3.51} (1.06) \\
Workload /100 & 62.46 (11.26) & 49.50 (12.93) & \textbf{46.06} (15.28) \\
\bottomrule
\end{tabular}
\end{table}

After correction across the 25 secondary tests, both assisted conditions reduced mental demand and increased perceived performance relative to control (all $p<.001$). Constant reduced effort ($p=.004$); adaptive's effort comparison was not significant ($p=.078$). Adaptive reduced frustration versus control ($p=.002$), whereas constant did not ($p=.086$). Physical and temporal demand comparisons were not significant. No workload dimension differed significantly between assisted policies.

Frustration averaged 4.69 in control, 3.93 in constant, and 3.51 in adaptive. Adaptive minus constant was $-0.42$ (CI [$-0.90$, 0.07], $t(23)=-1.77$, adjusted $p=.997$; Figure~\ref{fig:frustration}). This contrast did not establish lower frustration under adaptive than constant assistance.

\begin{figure}[!htbp]
\centering
\includegraphics[width=.94\columnwidth]{figures/frustration-tests.pdf}
\caption{Constant and adaptive frustration ($n=24$), mean $\pm$ SD. Lower 1--7 ratings indicate less frustration; paired-test adjusted $p=.997$.}
\label{fig:frustration}
\Description{Two bars compare constant frustration at3.93 and adaptive at3.51, with standard-deviation whiskers. The adjusted paired comparison is not significant.}
\end{figure}

\subsection{Intervention experience}
Adaptive received descriptively higher means on all four assistance items (Table~\ref{tab:assistance}). The timing contrast was $+0.56$ points ($t(23)=2.91$, raw $p=.008$), but not significant after secondary-family correction ($p=.119$). Helpfulness, seamlessness, and frustration-relief comparisons also did not pass correction (Figure~\ref{fig:assistance}). These ratings combine text hints and robot retrieval rather than isolating either modality.

\begin{table}[H]
\centering
\caption{Assistance ratings, mean (SD), $n=24$. All items use 1--7, with higher ratings indicating a better experience.}
\label{tab:assistance}
\begin{tabular}{lrr}
\toprule
Measure & Constant & Adaptive \\
\midrule
Helpfulness & 4.58 (1.26) & \textbf{4.83} (1.35) \\
Timing & 4.68 (1.48) & \textbf{5.24} (1.40) \\
Seamlessness & 4.01 (1.39) & \textbf{4.10} (1.39) \\
Frustration relief & 4.42 (1.35) & \textbf{4.82} (1.46) \\
\bottomrule
\end{tabular}
\end{table}

\begin{figure}[!htbp]
\centering
\includegraphics[width=\columnwidth]{figures/assistance-tests.pdf}
\caption{Constant and adaptive assistance ratings ($n=24$), mean $\pm$ SD. Higher 1--7 ratings indicate a better experience. Labels give adjusted paired-test $p$ values.}
\label{fig:assistance}
\Description{Four pairs of bars compare helpfulness, timing, seamlessness and frustration relief. Constant is blue and adaptive is orange. Adaptive means are descriptively higher; no adjusted paired comparison is significant.}
\end{figure}

\subsection{Overall experience}
Fifteen final surveys gave mean helpfulness 4.93 ($SD=1.62$), efficacy 4.87 ($SD=1.81$), trust 4.60 ($SD=1.76$), and agreement with following uncertain robot guidance 4.93 ($SD=1.67$).

Two final-survey comments illustrate the experience of assistance. One participant wrote, \textit{``The hints were helpful and overall well timed, but the actual robot arm moving the pieces were not adding anything more''}. Another wrote, \textit{``sometimes hints being given too often was counterintuitive and more distracting then better spaced hints, even though any hint helps.''} These comments distinguish the value of guidance from retrieval overhead and the costs of frequent intervention.

\section{Discussion, Limitations, and Future Work}

\subsection{Task progress and the experience of assistance}
Both assisted conditions support progress and reported workload relative to control. Constant has the highest observed completion and shortest mean duration; adaptive has better mean timing ratings and lower frustration. However, no corrected assisted-policy contrast establishes superiority. A practical choice between policies therefore remains unresolved, rather than demonstrating that adaptive assistance is preferable or equivalent.

Regular hints may provide repeated opportunities to advance an arrangement, whereas HR-rise-triggered hints may concentrate help at different moments. These are explanations to test, not mechanisms established by the present data. The final-survey comments suggest evaluating content and physical contribution separately: piece retrieval leaves placement to the participant and may add movement overhead, while frequent hints may interrupt reasoning.

\subsection{Physiological feedback and limitations}
The trigger detects an HR rise, not validated psychological stress. Wearable cardiac measurements are sensitive to movement and sensor contact \cite{quigley2024}. The observed assistance benefits therefore do not establish that the HR-rise flag measures stress or workload. Researchers observed some waiting for assistance, but did not systematically code its prevalence or condition distribution. Participants were not told the adaptive trigger rules. Future studies could test explanations and user-requested help, while accounting for PACE's finding that explicit queries can increase waiting \cite{delazzari2025}.

The sample consists of 24 college-recruited volunteers aged 20--27 years ($M=21.96$), mostly undergraduates with graduate students/researchers and faculty also participating. This age range and recruitment context limit generalization to older adults and experienced industrial workers. Balanced orders do not remove puzzle, period, device, or learning effects from the paired analysis. Frustration is a component of the workload composite, so the two measures are not independent.

Future work should isolate hint delivery from retrieval or direct placement, relate intervention events to observable progress, and use models that account for puzzle and order. Tangrams share orientation and placement demands with assembly \cite{zhao2025}, while retrieval resembles material provision in order picking \cite{smit2024}. Transfer to either setting requires its own safety, precision, and workflow evaluation.

\section{Ethical Considerations}

Participation was voluntary and consent was obtained before the task, with the option to pause or stop without penalty. Data comprised cardiac measurements, survey responses, and consented workspace recording/live viewing limited to the table and hands. Participant codes identify analysis records; any release must protect identifying information and respect consent coverage. Concealed Wizard-of-Oz operation was not disclosed during or after the task.

\section{Conclusion}

We compared unassisted, periodic, and HR-rise-triggered assistance in 24 participants' physical tangram attempts. Both assisted conditions yielded greater correct-piece progress, higher completion, and lower adapted workload than control. Constant had the best observed task means; adaptive had favorable intervention-experience means, without corrected evidence of an assisted-policy difference. The results support assistance in this task while leaving policy choice, physiological validity, and the added value of robot retrieval open.

\section*{AI Assistance Disclosure}

AI-assisted tools were used in preparing this manuscript. The authors are responsible for the accuracy, originality, and integrity of the work, including its citations.

\bibliographystyle{ACM-Reference-Format}
\bibliography{references}
\end{document}
